%==============================================================================
% tcc.tex --- Evaluating the energy cost of confidence-based offloading
% policies in hierarchical LLM inference cascades over passive optical
% networks
%
% Metadata, the abstracts (no text yet), the table of contents and lists, the
% chapter/section skeleton and the references. Running text so far: the
% Introduction's Context section; the rest is still to be written.
%
% Written in English (class option `english'); UFRGS also requires an
% abstract in Portuguese, the `translatedabstract' environment below.
%
% Build with:
%   pdflatex tcc
%   bibtex   tcc
%   pdflatex tcc
%   pdflatex tcc
%==============================================================================

\documentclass[cic,dipl,english]{infufrgs}

\usepackage[utf8]{inputenc}
\usepackage{graphicx}
\usepackage{amsmath}

%------------------------------------------------------------------------------
% Document metadata
%------------------------------------------------------------------------------
\title{Evaluating the energy cost of confidence-based offloading policies in
  hierarchical LLM inference cascades over passive optical networks}

% Title of the Portuguese abstract (Resumo)
\translatedtitle{Avalia\c{c}\~ao do custo energ\'etico de pol\'iticas de
  offloading baseadas em confian\c{c}a em cascatas hier\'arquicas de
  infer\^encia de LLMs sobre redes \'opticas passivas}

\author{Amorim}{Diego Hommerding}

\advisor[Prof.~Dr.]{Nazar}{Gabriel Luca}
\coadvisor[Prof.~Dr.]{Balreira}{Dennis Giovani}

%------------------------------------------------------------------------------
% Keywords
%------------------------------------------------------------------------------
\keyword{Distributed LLM inference}
\keyword{Passive optical networks}
\keyword{Energy efficiency}
\keyword{Confidence-based offloading}
\keyword{Fog computing}

% Keywords of the Portuguese abstract (Resumo)
\translatedkeyword{Infer\^encia distribu\'ida de LLMs}
\translatedkeyword{Redes \'opticas passivas}
\translatedkeyword{Efici\^encia energ\'etica}
\translatedkeyword{Offloading baseado em confian\c{c}a}
\translatedkeyword{Computa\c{c}\~ao em n\'evoa}

%==============================================================================
\begin{document}
%==============================================================================

\maketitle

%------------------------------------------------------------------------------
\begin{abstract}
\mbox{} % The class ends the abstract with a \\ (before the keywords), which
        % needs a line already open; this empty \mbox{} avoids a build
        % error until the abstract is written.
\end{abstract}

%------------------------------------------------------------------------------
\begin{translatedabstract}
\mbox{} % Portuguese abstract (Resumo); same placeholder as above.
\end{translatedabstract}

\tableofcontents

\listoffigures

\listoftables

%==============================================================================
\chapter{Introduction}
\label{chap:introduction}
%==============================================================================

As the use of Large Language Models (LLMs) becomes widespread in several domains of modern
society~\cite{liang2025adoption}, even if attitudes towards it still vary widely by age
and country~\cite{oecd2025technologies}, the challenge of adapting the network
architecture to handle this rapidly evolving technology remains an active research
topic~\cite{chen2026network}.
Since evolution seldom happens in isolation, the network itself has also been undergoing
changes, one of which is the adoption of Passive Optical Network (PON)
technology~\cite{effenberger2007introduction}, driven by the demand for greater capacity,
lower latency and reduced operating costs, alongside or in place of the hybrid
fibre-coaxial plant that DOCSIS runs over~\cite{prodnik2025migration}. Quality of service, scalability, security, and energy
cost are just some examples of the broad topics that need to be taken into account as
both of these changes happen in parallel, energy cost being the focus of this thesis.


LLM inference is expensive, and one way of making it cheaper is the model cascade: a
smaller model answers what it can, and the query is only passed to a larger one when its
answer is not reliable enough~\cite{li2021cascadebert,chen2024frugalgpt}. The idea is
older than the models it is applied to today, and what counts as expensive depends on the
setting, from the layers a classifier runs to the price of an API call.

Where the models run is a separate question. With edge computing, inference has started to
move out of the cloud and closer to the user, spread across end devices, edge nodes and
cloud servers, each tier with more computing power than the one below it and further
away~\cite{chen2026network}. RecServe puts the two ideas together, placing a cascade on
the tiers of such a network and letting each tier decide, from the confidence scores it
has seen recently, whether to answer the query or to pass it up~\cite{wu2025recserve}.

\section{Use of Artificial Intelligence Tools}
\label{sec:ai-tools}
Generative Artificial Intelligence tools were used during the preparation of this work,
both as a writing aid for parts of this text and as an assistant in its implementation.
All AI-assisted content was critically reviewed, edited, and validated by the author, who
assumes full responsibility for the content, analyses, and conclusions presented in this
work.

%==============================================================================
\chapter{Literature Review}
\label{chap:review}
%==============================================================================

\section{Recursive confidence-based offloading: RecServe}
\label{sec:recserve}

\section{Hierarchical inference over TWDM-PON}
\label{sec:pakpahan}

\section{Energy metrics for LLM inference}
\label{sec:metrics}

\section{Energy measurement by hardware class}
\label{sec:measurements}

\subsection{User device}

\subsection{Edge}

\subsection{Fog}

\subsection{Cloud}

\section{Challenges Addressed in This Work}
\label{sec:challenges}

%==============================================================================
\chapter{Methodology}
\label{chap:methodology}
%==============================================================================

%==============================================================================
\chapter{Results}
\label{chap:results}
%==============================================================================

%==============================================================================
\chapter{Conclusion}
\label{chap:conclusion}
%==============================================================================

%==============================================================================
% References
%==============================================================================
% ABNT author-date style with English labels (see README.md)
\bibliographystyle{abntex2-alf-en}
\bibliography{tcc}

\end{document}
